In order to achieve low hardware complexity, we lay out all the qubits on the first tier as a grid of $|G|$ $3 \times 3$ grids $\mathcal{G}_g$ as shown in the bottom center panel of Fig.~\ref{fig:thickness3_layout}. We refer to each $\mathcal{G}_g$ as a module. The edges in each module only correspond to the connections from the first group element in each ring entry of the base matrices used to construct the mitten code. To realize the full connectivity of the Tanner graph, we require additional inter-module connections between qubits in different $\mathcal{G}_g$.

\smallskip
\noindent\textbf{Placement as a Quadratic Assignment Problem.}
Each module $\mathcal{G}_g$ occupies one cell of a $\lceil\sqrt{|G|}\rceil \times \lceil\sqrt{|G|}\rceil$ super-grid on the chip. Since the intra-module (first-group-element) edges are absorbed into tier~0 by construction, the only free choice affecting the higher-tier edge lengths is the assignment $\pi : G \to \{0,\ldots,|G|-1\}$ that maps each module $\mathcal{G}_g$ to a slot on the super-grid. Given $\pi$, the total Manhattan length of the inter-module edges is
\begin{equation}
    C(\pi) \;=\; \sum_{(u,v)\in E_{1}\cup E_{2}} \bigl\lVert \mathrm{pos}_\pi(u) - \mathrm{pos}_\pi(v)\bigr\rVert_1,
    \label{eq:qap_cost}
\end{equation}
where $E_1,E_2$ are the layer-1 and layer-2 edge sets of the thickness-3 decomposition and $\mathrm{pos}_\pi(\cdot)$ is the chip coordinate implied by $\pi$. Minimizing $C(\pi)$ over all $|G|!$ permutations is an instance of the Quadratic Assignment Problem (QAP), which is NP-hard in general and empirically intractable for $|G|\gtrsim 40$ even with commercial branch-and-bound solvers~\cite{qap_survey}. Formulating (\ref{eq:qap_cost}) as a mixed-integer linear program and solving it with Gurobi on the $|G|=30$ instance for $20$ minutes leaves a $79\%$ LP relaxation gap, so the linear formulation cannot prove optimality at this scale.

\smallskip
\noindent\textbf{Heuristic optimization.}
We therefore approach $C(\pi)$ heuristically in two phases controlled by a single budget parameter $n_\mathrm{iter}$. The first phase is a multi-restart best-improvement local search. For each of $\lfloor n_\mathrm{iter}/2 \rfloor$ restarts, we sample a uniformly-random permutation and then repeatedly apply the pairwise swap $(i,j)$ that decreases $C(\pi)$ the most among all $\binom{|G|}{2}$ candidate swaps, until no improving swap exists (i.e. the current permutation is a local minimum under transpositions). Cost evaluations are vectorized so a full pass through the $\binom{|G|}{2}$ candidates takes about $10$~ms at $|G|=30$ and about $1$~s at $|G|=180$; a complete greedy descent to a local minimum needs $5$–$20$ such passes. Different random restarts land in different basins of attraction; we keep the best local optimum found across restarts.

The second phase is a simulated-annealing (SA) refinement of the single best permutation $\pi^\star_1$ produced by Phase~1. Starting from $\pi^\star_1$, we run $50 \cdot n_\mathrm{iter}$ Metropolis iterations. At each iteration we pick two indices $i,j$ uniformly at random, swap $\pi(i)\leftrightarrow\pi(j)$, compute the change $\Delta$ in $C(\pi)$, and accept the swap if $\Delta \le 0$ or with probability $\exp(-\Delta/T)$ if $\Delta>0$; otherwise the swap is reverted. The temperature $T$ decays exponentially from $T_0 = 50$ down to $T_{\min} = 0.5$ over the run, so the walker starts by freely exploring the neighbourhood of $\pi^\star_1$ and gradually contracts toward greedy descent. Unlike Phase~1 this phase can escape shallow local minima at the cost of occasional worsening moves; because the walker may drift into regions of higher cost, we track separately the lowest-cost permutation ever visited along the trajectory. The final output is whichever is smaller: $\pi^\star_1$ itself, or the best point visited by SA.

\smallskip
\noindent\textbf{Spectral placement.}
As an alternative to local search we tried a deterministic \emph{spectral} placement that turns the combinatorial layout problem into a linear-algebra one~\cite{spectral_layout}. Let $W_{ij}$ be the number of inter-module edges between $\mathcal{G}_i$ and $\mathcal{G}_j$, $D$ the diagonal degree matrix with $D_{ii}=\sum_j W_{ij}$, and $L = D - W$ the graph Laplacian of the module-connectivity graph. If we assign each module a continuous 2D position $x_i \in \mathbb{R}^2$, the sum of squared edge lengths weighted by connectivity is
\begin{equation}
    \sum_{i,j} W_{ij}\,\lVert x_i - x_j\rVert_2^2 \;=\; \sum_{d\in\{1,2\}} v^{(d)\top} L\, v^{(d)},
    \label{eq:spectral_obj}
\end{equation}
where $v^{(d)}$ is the vector of $d$-th coordinates. To avoid the trivial collapse $x_i = 0$ we require each $v^{(d)}$ to be unit-norm and orthogonal to the all-ones vector $\mathbf{1}$, and the two coordinate axes to be mutually orthogonal. Standard spectral theory then gives the optimum in closed form: $L$ has $\mathbf{1}$ as its zero-eigenvalue eigenvector (the excluded constant mode), the eigenvector $v_2$ of the second-smallest eigenvalue is the optimal 1D embedding (the Fiedler vector), and $v_3$ is the optimal direction orthogonal to $v_2$. Placing module $i$ at
\begin{equation}
    f_i \;=\; \bigl(v_2[i],\, v_3[i]\bigr)\;\in\;\mathbb{R}^2
    \label{eq:spectral_coord}
\end{equation}
therefore gives the provably optimal continuous 2D embedding under (\ref{eq:spectral_obj}). Computing the two eigenvectors requires one eigendecomposition of a $|G|\times|G|$ symmetric matrix and takes tens of milliseconds even for $|G|=180$.

The spectral coordinates $\{f_i\}$ are continuous, but the modules must occupy the $|G|$ integer slots of the super-grid, one module per slot. Finding the bijection $\pi$ that minimizes the total displacement $\sum_i \lVert f_i - g_{\pi(i)}\rVert_2^2$ (where $g_s$ is the coordinate of slot $s$) is the classical \emph{assignment problem}, which the Hungarian algorithm~\cite{hungarian} solves optimally in $O(|G|^3)$ time by primal-dual iteration on a bipartite matching between modules and slots. Concretely, one starts from a feasible dual (row and column potentials that keep every reduced cost non-negative), grows a matching using only zero-reduced-cost edges, and updates the potentials whenever no augmenting path exists; the procedure terminates at a globally optimal assignment in polynomial time. The full spectral pipeline\textemdash eigendecomposition plus Hungarian assignment\textemdash therefore delivers a globally optimal solution to the continuous relaxation of the placement QAP and its optimal snap-to-grid rounding in a single deterministic pass, with no random restarts and total wall time under one second for every code we tested.

Spectral placement produced the best hardware complexity we found for $[[300,60,14]]$ ($C_\mathrm{hw}=2.37$), beating $20$ multi-restart SA configurations that all landed near $C_\mathrm{hw}=2.46$. For all other codes in our family, however, multi-restart SA from random permutations dominated the spectral placement, presumably because the spectral objective (\ref{eq:spectral_obj}) minimizes the sum of \emph{squared} edge lengths, whereas HAL's routing cost is more closely tied to the sum of \emph{Manhattan} lengths (\ref{eq:qap_cost}) and to the discrete integer positions that SA can search over directly.

Because HAL's downstream $A^*$ routing objective (hardware complexity) is not monotone in $C(\pi)$, we treat (\ref{eq:qap_cost}) as a proxy and select the final layout by running HAL on each candidate permutation and keeping the one with the lowest hardware complexity, rather than the one with the lowest $C(\pi)$.

\smallskip
\noindent\textbf{HAL configuration.}
We swept a small number of HAL knobs to identify a configuration that consistently produced strong layouts across all codes tested: routing grid size $200 \times 200$, MPS edge order \texttt{length\_asc}, routing edge order \texttt{sum\_degrees\_desc}, at most $10$ face-switches per coupler, and (most impactfully) node and edge expansion set to $0$ so that each placed node and routed edge occupies only its own grid cell rather than blocking an eight-cell ring around it. With these settings HAL absorbs all first-group-element edges into tier~0 and packs the remaining $E_1 \cup E_2$ edges into $6$ to $17$ higher tiers depending on code size, giving the hardware complexities reported in Table~\ref{tab:paper_nonab_hw}.
